\documentclass[10pt,a3paper]{article}
\usepackage[margin=10mm]{geometry}
\usepackage[T1]{fontenc}
\usepackage[utf8]{inputenc}
\usepackage[french]{babel}
\usepackage{lmodern}
\usepackage{microtype}
\makeatletter\def\input@path{{../../../Style/}}\makeatother
\NeedsTeXFormat{LaTeX2e}
\ProvidesPackage{TLPosterCarteMentale}[2026/08/03 Posters et cartes mentales SI]

\RequirePackage{tikz}
\RequirePackage[most]{tcolorbox}
\RequirePackage{xcolor}
\RequirePackage{amsmath,amssymb}
\RequirePackage{graphicx}
\usetikzlibrary{positioning,arrows.meta,calc,shapes.geometric,mindmap,backgrounds,fit}

\definecolor{TLPosterBlue}{RGB}{24,86,141}
\definecolor{TLPosterLight}{RGB}{234,243,250}
\definecolor{TLPosterAccent}{RGB}{232,126,34}
\definecolor{TLPosterGreen}{RGB}{52,152,102}
\definecolor{TLPosterRed}{RGB}{192,57,43}
\definecolor{TLPosterGray}{RGB}{80,90,100}

\newcommand{\TLPosterTitre}[2]{%
  \begin{tcolorbox}[enhanced,boxrule=0pt,arc=0pt,colback=TLPosterBlue,left=6mm,right=6mm,top=4mm,bottom=4mm]
    {\color{white}\bfseries\LARGE #1}\hfill{\color{white}\large #2}
  \end{tcolorbox}%
}

\newtcolorbox{TLPosterBloc}[2][]{
  enhanced,breakable,
  colback=TLPosterLight,
  colframe=TLPosterBlue,
  colbacktitle=TLPosterBlue,
  coltitle=white,
  fonttitle=\bfseries,
  title={#2},
  boxrule=0.8pt,
  arc=2mm,
  left=3mm,right=3mm,top=2mm,bottom=2mm,
  #1
}

\newtcolorbox{TLPosterFormule}[1][]{
  enhanced,
  colback=white,
  colframe=TLPosterAccent,
  boxrule=1pt,
  arc=2mm,
  fontupper=\bfseries,
  halign=center,
  #1
}

\newcommand{\TLPosterMethode}[1]{%
  \begin{tcolorbox}[enhanced,colback=TLPosterGreen!8,colframe=TLPosterGreen,title=Méthode,fonttitle=\bfseries]
  #1
  \end{tcolorbox}%
}

\newcommand{\TLPosterAttention}[1]{%
  \begin{tcolorbox}[enhanced,colback=TLPosterRed!6,colframe=TLPosterRed,title=Point de vigilance,fonttitle=\bfseries]
  #1
  \end{tcolorbox}%
}

\tikzset{
  TLmindroot/.style={concept,concept color=TLPosterBlue,text=white,font=\bfseries\large,minimum size=34mm},
  TLmindlevel1/.style={level 1 concept/.append style={font=\bfseries,minimum size=23mm,sibling angle=45}},
  TLmindlevel2/.style={level 2 concept/.append style={font=\small,minimum size=17mm}},
  TLarrow/.style={-{Stealth[length=3mm]},thick,draw=TLPosterBlue},
  TLbox/.style={rounded corners=2mm,draw=TLPosterBlue,fill=TLPosterLight,align=center,inner sep=3mm}
}

\newenvironment{TLCarteMentale}[1]{%
  \begin{tikzpicture}[mindmap,grow cyclic,concept color=TLPosterBlue,every node/.style={concept},TLmindlevel1,TLmindlevel2]
  \node[TLmindroot]{#1}
}{;
  \end{tikzpicture}
}

\newcommand{\TLBranche}[3][]{child[concept color=#2]{node[#1]{#3}}}

\endinput

\pagestyle{empty}
\renewcommand{\familydefault}{\sfdefault}
\begin{document}
\TLPosterCourseHeader{13}{Machines asynchrones et synchrones}{Triphasé, onduleur et conversion électromécanique en alternatif}
\vspace{2mm}
\begin{minipage}[t]{0.61\linewidth}
\begin{TLPosterSavoirs}
\begin{itemize}
  \item maîtriser tensions simples/composées, couplages étoile/triangle et puissances triphasées
  \item comprendre la génération triphasée par onduleur et la modulation de largeur d'impulsion
  \item définir vitesse synchrone, glissement, bilan de puissance et couple de la MAS
  \item comprendre commande à $V/f$ constant et influence de la fréquence
  \item modéliser la MS avec diagramme de Behn--Eschenburg, couple électromagnétique et autopilotage brushless
\end{itemize}
\end{TLPosterSavoirs}
\end{minipage}\hfill
\begin{minipage}[t]{0.36\linewidth}
\TLPosterKey{Synchronisme}{$n_s=\dfrac{60f}{p}$\qquad $g=\dfrac{n_s-n}{n_s}$ pour la machine asynchrone}
\end{minipage}
\vfill
\begin{minipage}[t]{0.49\linewidth}
\begin{TLPosterBloc}[Système triphasé]{TLPosterBlue}
Trois tensions sinusoïdales déphasées de 120° alimentent les phases. Les couplages étoile et triangle relient tensions simples et composées. Les puissances active, réactive et apparente caractérisent la charge.
\end{TLPosterBloc}
\begin{TLPosterBloc}[Machine asynchrone]{TLPosterPurple}
Le champ statorique tourne à la vitesse de synchronisme. Le rotor présente un glissement nécessaire à l'induction des courants et au couple. Le bilan de puissance distingue pertes statoriques, rotorique et mécanique.
\end{TLPosterBloc}
\end{minipage}\hfill
\begin{minipage}[t]{0.49\linewidth}
\begin{TLPosterBloc}[Onduleur triphasé]{TLPosterGreen}
À partir d'une source continue, six interrupteurs synthétisent des tensions triphasées. La MLI commande fréquence et amplitude du fondamental et donc les conditions d'alimentation de la machine.
\end{TLPosterBloc}
\begin{TLPosterBloc}[Machine synchrone]{TLPosterTeal}
Le rotor suit le champ tournant sans glissement en régime établi. Le diagramme de Behn--Eschenburg relie tension, courant et fcem. L'autopilotage synchronise la commande électronique avec la position rotorique.
\end{TLPosterBloc}
\end{minipage}
\vfill
\begin{TLPosterMethode}
\begin{enumerate}
  \item Identifier couplage, tensions, courants et nombre de paires de pôles.
  \item Calculer vitesse de synchronisme et, pour la MAS, le glissement.
  \item Établir le bilan des puissances et le couple utile.
  \item Exploiter le schéma ou diagramme vectoriel adapté à la machine.
  \item Relier fréquence et amplitude de l'onduleur au point de fonctionnement.
  \item Vérifier limites thermiques, couple, vitesse et rendement.
\end{enumerate}
\end{TLPosterMethode}
\vfill
\TLPosterRetenir{les machines alternatives se comprennent à partir du champ tournant, du synchronisme et du bilan de puissance.}
\end{document}
